\documentclass[11pt,a4paper]{article}
\usepackage[T1]{fontenc}
\usepackage{isabelle,isabellesym}
\usepackage{amsmath,amssymb}
\usepackage{pdfsetup}
% a wider text block: the deep embedding's statements and proofs need the line length
\usepackage[a4paper,textwidth=460pt,textheight=650pt]{geometry}
\urlstyle{rm}\isabellestyle{it}
% suppress "theory ... imports ... begin" headers and trailing "end"s
\isadroptag{theory}
\hyphenation{Benz-m\"ul-ler Kohl-hase Kirch-ner Isa-belle Paul-son}
\emergencystretch=5em % let TeX stretch otherwise-overfull lines rather than overrun the margin
\begin{document}
\title{HOL in HOL over a set-theoretic universe:\\
  consistency of $\mathrm{NK}$ with infinity and choice}
\author{Christoph Benzm\"uller \and Daniel Kirchner}
\maketitle
\begin{abstract}
This companion to the entry \emph{A deep embedding of HOL in HOL}
(\texttt{HOL\_in\_HOL\_Deep})~\cite{HOLinHOLDeep} certifies the
natural-deduction calculus $\mathrm{NK}$ of that entry consistent with a
genuine Dedekind-style \emph{axiom} of infinity---an injective,
non-surjective self-map of the individuals---and, because its standard model
has full function spaces, with the relational axiom of \emph{choice} as
well.  Such an axiom has \emph{no} finite models, so the compactness
argument of the main entry---which lifts arbitrarily large \emph{finite}
models to an infinity \emph{scheme}---does not apply: a genuinely
\emph{infinite} model is required.  We construct one over Paulson's
\texttt{ZFC\_in\_HOL}~\cite{ZFCinHOL}: the individuals are interpreted as the
finite ordinals $\omega$, with full function spaces at every type, yielding a
\emph{standard} $\Sigma$-model of $\mathrm{NK}$ over the set-theoretic
universe $V$.  As this meta-theory is stronger and non-conservative, the
guarantee is \emph{relative} consistency at $\mathrm{ZF}$ strength; the
machine-checked formalisation lives in Isabelle/HOL with Paulson's
\texttt{ZFC\_in\_HOL} (the $\mathrm{ZF}$ reading is metatheoretic, made
precise in the introduction).  Because $\mathrm{NK}$ already provides
extensionality and typed description (after Andrews), this certifies the
consistency of \emph{full} classical higher-order logic in Church's
sense, choice included as a separate relational scheme.

The infinity axiom and the weaker infinity \emph{scheme} of the main
entry turn out to be logically independent.  That the scheme does not
entail the axiom is shown in the main entry; the present entry settles
the other direction, since the axiom does not entail any of the scheme's
individual inequalities: its set-theoretic model interprets all
individual constants alike, so those inequalities fail while the axiom
holds.  Neither sentence, then, implies the other.

The plain-HOL results---soundness, Henkin completeness, and
finite-model consistency, including the compactness route to the
infinity scheme---are proved in the main entry~\cite{HOLinHOLDeep} and
imported here \emph{unchanged}.  What the present entry \emph{adds} is
the set-theoretic route to axioms that plain HOL alone cannot model.
\end{abstract}

\medskip\noindent
This contribution is part of the LogiKEy project
(\url{https://logikey.org}).

\paragraph{Use of generative AI.}
Parts of this formalisation and its documentation were drafted with a
generative AI assistant (Anthropic's Claude) under continuous human
direction.  All proofs are machine-checked by Isabelle; the entire
development, including all comments and references, was reviewed, revised
and simplified by the authors, who take full responsibility for its content.

\setcounter{tocdepth}{1}
\tableofcontents
\newpage

\section{Introduction}

\paragraph{Two entries, and where each result is proved.}
This is the second of two Archive of Formal Proofs entries on a deep
embedding of classical higher-order logic---Church's simple theory of
types, following Benzm\"uller, Brown and Kohlhase~\cite{BBK04}---inside
Isabelle/HOL.  The first entry, \emph{A deep embedding of HOL in HOL}
(\texttt{HOL\_in\_HOL\_Deep})~\cite{HOLinHOLDeep}, stays \emph{entirely within
plain Isabelle/HOL, with no set-theoretic universe}.  There the
natural-deduction calculus $\mathrm{NK}$ is proved sound, Henkin-complete and
consistent; consistency comes from a concrete finite standard model, and, by
compactness over arbitrarily large \emph{finite} models, $\mathrm{NK}$ is
shown consistent even with an infinity \emph{scheme} $\{c_i \neq c_j \mid i
\neq j\}$---still without ever constructing an infinite model.  The present
entry adds one result that plain HOL cannot reach: the consistency of
$\mathrm{NK}$ with a genuine Dedekind-style \emph{axiom} of infinity, and
with the axiom of \emph{choice}.  It imports the first entry unchanged and
supplies the set-theoretic model that this result needs.

We are careful throughout---here, in the theory comments, and in both
abstracts---to say which of two tiers a statement belongs to: what holds in
plain HOL, and what needs the set-theoretic universe.  Within this entry the
split is visible in the theory structure itself.  The set-theory-free half,
theory \texttt{HOL\_Universe}, imports only the first entry and never
\texttt{ZFC\_in\_HOL}.  It contains the choice scheme with its satisfaction
condition (the infinity axiom and its satisfaction are inherited from the
first entry), the abstract relative-consistency theorem, the completeness
results, and two universe interfaces: the locale \texttt{hol\_universe}, and
the (at most as strong) Harrison-style axiom \texttt{harrison\_universe},
after Harrison's self-verification of HOL Light~\cite{Harrison06}, from which
\texttt{hol\_universe} is in turn \emph{constructed}.  The set-theoretic
universe itself appears only in the second theory, which instantiates both
interfaces with Paulson's set-theoretic universe~$V$ (the type of all sets of
\texttt{ZFC\_in\_HOL}, introduced below).

\paragraph{Why the meta-theory must be strengthened.}
A single Dedekind-style axiom of infinity has \emph{no} finite models, by the
pigeonhole principle.  So the route of the first entry---compactness over
finite models---yields nothing here, and a genuinely \emph{infinite} model is
required.  We do not try to build one inside plain HOL, for two reasons.
First, no type that plain HOL can construct is large enough to carry a
universe of full function spaces: by Cantor's theorem each such type is
outgrown by the tower of function spaces built over it, and that no other
plain-HOL route succeeds follows from G\"odel's second theorem (discussed
below).  Second, the term-model construction of the completeness proof would
be circular for a single sentence.  Instead we borrow a ready-made, stronger
meta-theory: Paulson's \texttt{ZFC\_in\_HOL}~\cite{ZFCinHOL}, which
axiomatises a Zermelo--Fraenkel universe as a type $V$.  This is \emph{not} a
conservative extension of HOL---$V$ is postulated by axioms, and $\mathrm{ZF}$
is proof-theoretically stronger than plain HOL---so what we obtain is a
\emph{relative} consistency result: $\mathrm{NK}$ with infinity is consistent
\emph{relative to} $\mathrm{ZF}$.  (The formalisation rests on the full
\texttt{ZFC\_in\_HOL}, whose set-theoretic choice axiom we never invoke
essentially; that the
construction in fact needs only $\mathrm{ZF}$ strength---indeed less---is made
precise in \emph{How little of $\mathrm{ZF}$ is used} below.  This is why we
write $\mathrm{ZF}$, not $\mathrm{ZFC}$.)  This pins down exactly where set theory is
used, namely only to certify consistency (an axiom without finite models has
no finite consistency witness).  The model that realises the consistency is,
by contrast, a \emph{countable} Henkin model built in plain HOL, by the
(ZFC-free) model-existence half of the first entry.  For that single axiom
the stronger meta-theory supplies the consistency premise and nothing else.

\paragraph{Choice, and full HOL.}
The same standard model settles more.  Because its function spaces are
\emph{full}, it validates every instance of the relational axiom of
\emph{choice}: the required choice function is simply the set-theoretic
$\lambda$-abstraction of a meta-level Hilbert choice.  Hence $\mathrm{NK}$
with the Dedekind axiom of infinity \emph{and} choice is consistent relative
to $\mathrm{ZF}$.  Since $\mathrm{NK}$ already carries extensionality and
typed description, this is the consistency of full classical higher-order
logic in the sense of Church.  To be precise, by \emph{full classical HOL in
the sense of Church} we mean exactly this system: Church's monomorphic simple
theory of types~\cite{Church40}---the calculus $\mathrm{NK}$ with Boolean and
functional extensionality and typed description, extended by the Dedekind
axiom of infinity and the relational axiom of choice.  It is \emph{not} the
polymorphic Isabelle/HOL of the meta-level, which goes beyond Church's system
with type variables, type definitions and overloading.

\paragraph{Two tiers of consistency, and G\"odel's second theorem.}
Write \emph{O-HOL} for the embedded object logic (whose calculus is
$\mathrm{NK}$) and \emph{M-HOL} for the meta-logic Isabelle/HOL in which it
lives; the two must be kept apart.  The first entry proves
$\text{M-HOL}\vdash\mathrm{con}(\text{O-HOL})$: the bare object logic is
consistent, and this is shown in plain HOL with \emph{no} set universe.  (With
a single individual every domain is finite, so M-HOL exhibits a model
directly; and O-HOL, having finite models, lacks the arithmetic strength that
G\"odel's second theorem presupposes.)  The present entry adds the
\emph{full} logic:
$\text{M-HOL}+V\vdash\mathrm{con}(\text{O-HOL}+\text{infinity}+\text{choice})$,
that is, M-HOL extended by the set universe $V$ proves full classical HOL
consistent.  This is as close as one can come to a self-standing
$\vdash\mathrm{con}(\text{HOL})$, and it is emphatically \emph{not} the object
logic proving its own consistency.

That the meta-logic here cannot be plain M-HOL either is a
\emph{metatheoretic} observation, outside the machine-checked development.
With infinity and choice, O-HOL is itself Church's simple type theory with
infinity and choice, and this interprets the M-HOL core (its individuals
becoming an infinite type, its function types the function types).  So the two
share their proof-theoretic strength: by the classical analysis both are
equiconsistent with a Zermelo-style set theory~\cite{Mathias01}.  (Here
``M-HOL'' means the HOL \emph{core}, whose definitional and overloading
apparatus is separately known consistent~\cite{KuncarPopescu19}.)  Were plain
M-HOL to prove $\mathrm{con}(\text{O-HOL}+\text{infinity}+\text{choice})$, it
would thereby prove its own consistency, which G\"odel's second
incompleteness theorem forbids; so the step up to the set universe $V$ is
\emph{necessary}, not a convenience.  The same obstacle, and the same step to
a stronger meta-theory, appears in Harrison's self-verification of HOL
Light~\cite{Harrison06} and its successors~\cite{Kumar16}.  Throughout it is
M-HOL reasoning \emph{about} O-HOL, one level up---a difference of role, not
of logical power: the object logic is in no way the weaker system.  The
certificate is tied to exactly the axioms modelled here.  Extending O-HOL
with a further axiom stays legitimate, but the machine-checked guarantee does
not extend automatically: one must exhibit a model of the new axiom, and if
the extension reaches the consistency strength of $\text{M-HOL}+V$, certifying
it needs a yet stronger meta-theory---again by G\"odel's second theorem.

\paragraph{How little of $\mathrm{ZF}$ is used.}
A closer look pins down the exact surplus over plain HOL.  The object-side
argument is carrier-generic and uses no set theory at all: full HOL is
consistent as soon as \emph{some} standard $\Sigma$-frame has
Dedekind-infinite individuals (an injective, non-surjective self-map).  Plain
HOL has infinity already; what it lacks is a single type gathering the full
function spaces over an infinite base (a model formalised in HOL keeps all its
domains inside one carrier type)---by Cantor these outgrow every type it
can construct---and that is the whole reason for a stronger universe.  Choice
is \emph{not} part of this surplus: the same full model validates it for free,
so the price of full HOL over plain HOL is the axiom of infinity alone.  Nor
need that universe be full $\mathrm{ZF}$.  It is enough to assume full
function spaces, booleans, extensionality, infinitely many individuals and a
domain-respecting parameter interpretation---no replacement and no
ordinals---and from these the logical constants (negation, disjunction,
quantification, equality and description) are \emph{constructed} rather than
assumed.  Every domain then has bounded rank, and this is machine-checked:
the domain of individuals sits at rank $\omega$, and each function-space layer
adds only finitely much, so $\mathrm{rank}(D_\sigma) < \omega+\omega$ for
every type $\sigma$, and every domain is an element of the \emph{set}
$V_{\omega\cdot 2}$.  Two things should be kept apart here.  What is
\emph{machine-checked} is precisely this rank bound together with
$\text{Isabelle/HOL}+\texttt{ZFC\_in\_HOL}\vdash
\mathrm{con}(\text{O-HOL}+\text{infinity}+\text{choice})$.  That the same
consistency follows from far less than $\mathrm{ZF}$ is a
\emph{metatheoretic} reading of that proof, read off from the ranks of the
objects it constructs: a single derivation mentions only finitely many
types, so its refutation is validated using only sets of bounded
rank---comfortably within Zermelo set theory, and without ever forming
$V_{\omega\cdot 2}$ as a completed set (which would already call on a trace
of replacement).  This reading is not itself a formalised
relative-consistency theorem.  Under it the relative consistency factors as
$\mathrm{con}(\mathrm{ZF}) \Rightarrow \mathrm{con}(\text{weak universe})
\Rightarrow \mathrm{con}(\text{full HOL})$; the first implication also
absorbs the choice that the meta-theory uses throughout---by G\"odel's
constructible universe $L$ at $\mathrm{ZF}$ strength, while at Zermelo
strength the corresponding elimination of choice without replacement is
genuinely delicate (see Mathias~\cite{Mathias01}).  The surplus strength of
\texttt{ZFC\_in\_HOL}---which \emph{exceeds} $\mathrm{ZF}$, as it proves
$\mathrm{con}(\mathrm{ZFC})$---is never used.

\input{session}
\bibliographystyle{abbrv}
\bibliography{root}
\end{document}
